% Standalone document: no external bibliography, figures, or input files.
% Compile with: pdflatex phase1_evidence_assessment.tex (twice for references).
\documentclass[11pt,letterpaper]{article}
\usepackage[margin=1in]{geometry}
\usepackage[T1]{fontenc}
\usepackage{lmodern}
\usepackage{microtype}
\usepackage{amsmath}
\usepackage{booktabs,tabularx,array,longtable}
\usepackage{float}
\usepackage{xcolor}
\usepackage{enumitem}
\usepackage[hidelinks]{hyperref}
\usepackage{xurl}

\definecolor{supported}{RGB}{24,96,65}
\definecolor{conditional}{RGB}{130,83,0}
\definecolor{missing}{RGB}{154,38,38}
\newcommand{\yes}{\textcolor{supported}{\textbf{Demonstrated}}}
\newcommand{\partiallydemonstrated}{\textcolor{conditional}{\textbf{Partially demonstrated}}}
\newcommand{\needed}{\textcolor{missing}{\textbf{Still needed}}}
\newcommand{\evidence}[1]{\textnormal{[E#1]}}
\newcolumntype{L}[1]{>{\raggedright\arraybackslash}p{#1}}
\newcolumntype{Y}{>{\raggedright\arraybackslash}X}
\setlength{\parindent}{0pt}
\setlength{\parskip}{6pt}
\clubpenalty=10000
\widowpenalty=10000
\brokenpenalty=10000
\setlist[itemize]{leftmargin=*,topsep=3pt,itemsep=3pt,parsep=0pt}
\setlist[enumerate]{leftmargin=*,topsep=3pt,itemsep=4pt,parsep=0pt}
\renewcommand{\arraystretch}{1.16}
\hypersetup{
  pdftitle={Phase I Evidence Assessment: Multi-Agent LLM Auction Experiments},
  pdfsubject={Demonstrated capabilities, existing results, and remaining evidence gaps}
}

\title{Phase I Evidence Assessment\\\large Multi-Agent LLM Auction Experiments}
\author{}
\date{October 6, 2026}

\begin{document}
\begin{center}
{\LARGE\bfseries Phase I Evidence Assessment\par}
{\large Multi-Agent LLM Auction Experiments\par}
October 6, 2026
\end{center}

\section{Assessment and scope}

\textbf{The existing work demonstrates a functioning multi-agent auction
environment, substantive market mechanics, behavioral classification and
comparison, and greater than 90\% allocative surplus efficiency in specified
recorded experiments. It does not yet establish every Phase I requirement.}
The principal outstanding items are verified operation with at least ten LLMs
and confirmation that the implemented public-information updates meet the
intended ``news feed'' requirement. A requirement for a trained behavioral
classifier would also need additional evidence.

This assessment applies the six Phase I requirements supplied for the related
DARPA SBIR effort to the repository accompanying \emph{LLM Bidders Preserve the
Mechanism-Level Orderings of Human Bidders}. It distinguishes executed
experiments from configurable capabilities and distinguishes previously
reported findings from a new analysis of already collected data. It does not
claim that these experiments were performed under, or funded by, the SBIR
grant. Collaboration, deception, influence, and development of human-like
reference agents are outside this Phase I assessment.

The evidence consists of the manuscript, simulator code, archived
configurations and prompts, recorded auction outcomes, and analysis outputs.
An offline audit calculated allocation efficiency from existing outcomes;
it did not generate new model responses. Evidence identifiers below refer to
the repository artifacts listed in Section~\ref{sec:artifacts}.

\begin{table}[H]
\small
\caption{Requirement-by-requirement assessment. Status applies to the scope stated.}
\label{tab:assessment}
\begin{tabularx}{\textwidth}{@{}L{0.25\textwidth}Y@{}}
\toprule
Requirement & Evidence status and remaining issue \\
\midrule
1. Interacting AI agents & \yes. Agents compete in executed auctions with shared allocation and payment rules. \evidence{1--3} \\
2. Market mechanics & \partiallydemonstrated. Multiple formats, bids, prices, allocations, simulated payoffs, timing, and changing public market information are implemented. Quantity is one object; an external news feed is not established. \evidence{2--4} \\
3. At least ten LLMs & \needed. Six distinct models have recorded evidence. A configurable EDSL interface is present, but ten-model execution has not been verified. \evidence{3,5} \\
4. Behavioral classifiers & \yes\ for rule-based economic classification and cross-model measurement. A trained classifier and its predictive validation are not established. \evidence{6--7} \\
5. Economic behavior; efficiency above 90\% & \yes\ for completed GPT-4o private-value baseline auctions under the surplus definition in Section~\ref{sec:efficiency}. This is not a universal result across models. \evidence{8} \\
6. AI--AI and AI--human comparison & \yes\ across model-specific experiments and published human benchmarks. Mixed-model auctions and participant-level human records are not established by this evidence. \evidence{1--2,7} \\
\bottomrule
\end{tabularx}
\end{table}

\section{What the existing experiments demonstrate}

\subsection{Multiple agents participate in an actual auction}

\yes. The simulator instantiates separate bidders with private values or
signals, elicits their actions through EDSL, and applies a common auction
mechanism to determine the winner and payments. The baseline private-value
experiments audited here have three bidders. Submitted bids affect other
agents' winning chances and prices; clock and proxy auctions additionally
support responses to an evolving auction state. This is an executed strategic
interaction, rather than a collection of independent hypothetical bidding
questions. Raw results contain bidder values, actions, a recorded winner,
price, profits, and available free-text plans. \evidence{2--3}

\textbf{Scope.} The main model comparisons use separate auctions populated by
one model at a time. Interaction among multiple agents is demonstrated;
competition among different model types within the same auction is not
demonstrated by the main experimental panel. \evidence{1}

\subsection{Market structures, payoffs, timing, and public information}

\yes\ for the following mechanics:
\begin{itemize}
  \item \textbf{Multiple structures and information environments.} The main
  paper covers first-price and second-price independent-private-value auctions,
  sealed second-price affiliated-private-value auctions, open and closed
  ascending-clock affiliated-private-value auctions, and first-price and
  second-price common-value auctions. Additional recorded experiments cover
  eBay-style proxy bidding and hard versus soft closing rules.
  \item \textbf{Bids, prices, assets, quantities, and timing.} Bidders compete
  for one object per auction. The simulator records numeric bids or dynamic
  bidding decisions, determines prices and allocations, and implements clock
  periods and proxy-auction closing rules. Asset quantity is fixed at one;
  variable-quantity trading and multi-asset portfolios are not evidenced here.
  \item \textbf{Incentives and utility.} Prompts instruct bidders to maximize
  profits. Allocation and payment rules produce simulated monetary payoffs;
  for the audited private-value auctions, the winner's payoff is value minus
  price and losers receive zero. These are numerical economic incentives,
  not actual cash transfers to the models.
  \item \textbf{Outcome measurement.} Existing analyses measure bidding
  deviations, profit deviations in common-value auctions, and timing. The
  recorded allocations and values also support the efficiency calculation
  in Section~\ref{sec:efficiency}.
\end{itemize}
These capabilities are supported by the simulator and archived runs, not just
by proposed configurations. \evidence{2--4,7--8}

\partiallydemonstrated\ for the ``news feed'' requirement. Open-clock bidders receive
updates on prices and dropouts; eBay bidders receive the current price and a
price-history transcript; repeated-auction experiments provide outcome
feedback across auctions. These constitute changing public \emph{market}
information. No separate stream of external news, public fundamental-value
shocks, or announcements is established by the inspected evidence.
\textbf{Still needed:} determine whether endogenous market updates satisfy
the requirement. If exogenous news is required, provide evidence of such a
feed and agents' responses to it, or implement and demonstrate that capability.
\evidence{2--4}

\subsection{Model coverage}

\partiallydemonstrated. The recorded data establish six distinct model identifiers:

\begin{center}
\small
\begin{tabular}{@{}ll@{}}
\toprule
Model & Recorded identifier \\
\midrule
GPT-4o & \texttt{gpt-4o} \\
Claude 3.5 Haiku & \texttt{claude-3-5-haiku-20241022} \\
Gemini 2.0 Flash & \texttt{gemini-2.0-flash} \\
GPT-5-mini & \texttt{gpt-5-mini} \\
Llama-3-8B-Instruct & \texttt{meta-llama/Meta-Llama-3-8B-Instruct} \\
Gemma-3-27B & \texttt{google/gemma-3-27b-it} \\
\bottomrule
\end{tabular}
\end{center}

The first five appear in the main model comparison; Gemma supplies additional
framing evidence. The historical directory suffix \texttt{claude\_sonnet}
contains Claude 3.5 \emph{Haiku} observations and must not be counted as a
separate model. Different temperatures or prompts likewise do not increase
the distinct-model count. \evidence{2,5}

The simulator accepts configurable model and provider identifiers through
EDSL. That interface provides a route to broader coverage, but is not a
demonstration of successful auctions with ten models. \textbf{Still needed:}
locate qualifying existing runs for at least four additional distinct LLMs,
or execute and archive them in the required environment. Preserve model
identifiers, configurations, valid actions, completed outcomes, and failures.
The current data also have missing or thin model--format cells, including
missing GPT-5-mini clock outcomes. \evidence{2--3,5}

\subsection{Behavioral classification and differences}

\yes\ for explicit rule-based classifiers. The analysis defines an economic
benchmark $b^*(v)$ and labels each bid as within a tolerance band, below the
benchmark, or above it:
\[
\begin{aligned}
\text{Within band:}\quad & |b-b^*(v)|\leq 0.02\,b^*(v),\\
\text{Underbid:}\quad & b<b^*(v)\quad\text{outside the band},\\
\text{Overbid:}\quad & b>b^*(v)\quad\text{outside the band}.
\end{aligned}
\]
For three-bidder first-price independent-private-value auctions the benchmark
is $b^*(v)=2v/3$; in the truthful formats it is $b^*(v)=v$.
The analysis reports the within-band share among all bids and the under/over
split among deviating bids. Clock bid-direction analyses exclude winners
because their exit values are censored. \evidence{6}

Existing outputs also report scaled mean absolute deviations, uncertainty
intervals, cross-format rankings, and framing effects. For example, first-price
private-value bid deviations are approximately 27.6\% for GPT-4o, 24.1\% for
Claude 3.5 Haiku, 27.1\% for Gemini, 0.1\% for GPT-5-mini, and 38.2\% for
Llama. These percentages describe the paper's normalized bidding-error metric,
\emph{not} allocation efficiency. \evidence{7}

There is also a measured timing response to market rules: the eventual
winner's last maximum-bid increase occurs in period 10 in 63.3\% of hard-close
auctions versus 6.7\% of soft-close auctions, with 30 recorded auctions per
treatment. \evidence{4}

\textbf{Still needed, if required by the classifier specification:} a trained
detector with defined labels, held-out evaluation, and predictive-performance
metrics. The existing rule-based classifications support statements about
observable bidding tendencies and differences; they do not by themselves
identify the psychological causes of those differences.

\subsection{Comparisons across models and with human behavior}

\yes. The main analysis compares five LLMs across seven auction settings,
with missing or thin cells identified. It compares deviation magnitudes,
directions, and mechanism-level orderings with human experimental benchmarks.
The three non-reasoning large models reproduce important human orderings,
including larger first-price than second-price deviations, while differing
from humans in error magnitudes and often in error directions.
\evidence{1--2,7}

\textbf{Qualification.} The human evidence consists of published summary
statistics and synthetic reconstructions calibrated to those statistics,
not observed participant-level human records. Some benchmark designs differ
in bidder count, and human category definitions need not match the LLM
tolerance rule exactly. This supports comparisons with \emph{comparable}
economic environments and specified normalization, not a claim of identical
protocols or validated human-level fidelity. \textbf{Still needed only if
explicitly required:} matched participant-level human data, or auctions in
which different model types compete directly. \evidence{1--2}

\section{Allocative efficiency from existing data}
\label{sec:efficiency}

\yes\ for the recorded GPT-4o private-value baseline auctions under the
following explicit measure. For auction $a$, let $v_{ia}$ be bidder $i$'s
value and $w(a)$ the recorded winner. With one object, no reserve, and seller
opportunity cost taken as zero, realized total surplus is $v_{w(a),a}$ and
the maximum feasible surplus is $\max_i v_{ia}$. Payments transfer surplus
between buyer and seller. Aggregate allocative surplus efficiency is
\begin{equation}
\label{eq:efficiency}
E_{\mathrm{surplus}}=
100\times\frac{\sum_a v_{w(a),a}}{\sum_a\max_i v_{ia}}.
\end{equation}
This is distinct from both the unweighted mean of per-auction surplus ratios
and the frequency with which a highest-value bidder wins. The audit reports
all three. The main results below use Equation~\eqref{eq:efficiency}.

\begin{table}[H]
\centering
\small
\caption{GPT-4o allocation outcomes calculated from saved completed auctions.}
\label{tab:efficiency}
\begin{tabularx}{\textwidth}{@{}Yrrr@{}}
\toprule
Auction format & Auctions & \shortstack{Surplus\\efficiency} & \shortstack{Highest-value\\winner} \\
\midrule
First-price, independent private values & 95 & 99.77\% & 95.79\% \\
Second-price, independent private values & 90 & 99.72\% & 96.67\% \\
Second-price, affiliated private values & 97 & 99.09\% & 89.69\% \\
Open ascending clock, affiliated private values & 53 & 100.00\% & 100.00\% \\
Closed ascending clock, affiliated private values & 57 & 100.00\% & 100.00\% \\
\midrule
Total completed auctions & 392 & & \\
\bottomrule
\end{tabularx}
\end{table}

Each of the five GPT-4o settings exceeds 90\% surplus efficiency. The
unweighted mean per-auction efficiencies also exceed 90\% in all five settings
(99.42\%, 99.51\%, 99.08\%, 100\%, and 100\%, respectively).
These are new calculations from existing data, rather than efficiency
statistics previously reported in the manuscript. \evidence{8}

\textbf{Limits of the claim.} These descriptive results concern completed,
saved auctions and do not establish performance across failed or missing runs,
unseen auctions, or all models. Llama's second-price independent-private-value
cell has only 85.32\% surplus efficiency. High allocation efficiency also
does not imply equilibrium bidding, accurate human simulation, or uniformly
positive bidder profits.

If DARPA instead defines efficiency as the percentage of auctions awarded to
a highest-value bidder, the affiliated sealed second-price GPT-4o cell is
89.69\% and does not meet the threshold. The other four GPT-4o cells exceed
90\% on that measure. Value ties count as efficient allocations.
\textbf{Still needed:} align the acceptance metric and its scope---for example,
whether it must hold for every model and format---with the grant requirement.

\subsection{Audit and reproducibility}

The offline audit covers 802 completed raw auctions in five private-value
formats across the five main models. It reads recorded winners directly and
maps bidder names to the saved value arrays; it does not infer clock winners
from potentially tied exit prices. It checks model identifiers and relevant
configuration fields, records source-file hashes, and checks profit accounting.
It excludes common-value, eBay, framing, temperature, repeated-round, and
duplicated robustness datasets from the efficiency calculation. \evidence{8}

An independent comparison against existing bidder-level CSV exports matched
800 auction allocations and prices, with prices compared to eight decimal
places. Two Llama clock records had no matching CSV export. One Gemini open-clock
record reports winner profit 10 although recorded value 37 minus price 28
equals 9. Those discrepancies are preserved in the audit; the source data
were not altered. All 392 GPT-4o auctions in Table~\ref{tab:efficiency} match
their CSV exports and pass the profit-accounting check.

\textbf{Still needed for a clean broader evidence package:} reconcile the
Gemini profit discrepancy and document or restore the two missing Llama
exports. These issues do not affect the GPT-4o efficiency claim. An audit flag
for fewer than 30 completed auctions identifies small cells; it is not an
asserted DARPA sample-size requirement.

\begin{samepage}
To reproduce the audit from the repository root:
\begin{quote}
\ttfamily\small python3 scripts/audit\_phase1\_evidence.py
\end{quote}
\end{samepage}
The script uses the Python standard library and existing local records.
Outputs include a model inventory, per-auction efficiency records, a summary
table, and a JSON audit manifest. Export row counts in the model inventory
may overlap and must not be treated as independent observation counts.

\section{Remaining evidence needed}

\begin{enumerate}
  \item \textbf{Ten-model execution: outstanding demonstration.} Supply
  completed auction evidence for at least four additional distinct models,
  from existing archives if available or from additional runs. General EDSL
  availability is not sufficient evidence of execution in this environment.
  \item \textbf{News feed: resolve the requirement definition.} If changing
  endogenous public market information qualifies, document the existing
  price, dropout, and history updates. If external news is required,
  demonstrate an exogenous information feed and resulting agent behavior.
  \item \textbf{Classifier: resolve the required type and validation.} The
  current rule-based labels and comparisons are implemented and supported by
  outputs. A requirement for learned classification, predictive validation,
  or a broader taxonomy of biases requires additional evidence.
  \item \textbf{Efficiency: specify metric and coverage.} Use the defined
  surplus measure and the completed-run scope when claiming the threshold.
  A requirement that every model or format exceed 90\%, or that the result
  satisfy a statistical confidence bound, is not established by this audit.
  \item \textbf{Comparison scope and mechanics: confirm any stricter terms.}
  Variable quantities, multiple assets, direct mixed-model competition, and
  participant-level matched human experiments are not demonstrated here.
  They need further evidence only if the actual Phase I specification
  requires them.
  \item \textbf{Data reconciliation: limited cleanup.} Resolve the one
  recorded profit discrepancy and the two missing CSV exports before
  presenting the full cross-model audit as reconciled.
\end{enumerate}

These items concern evidence and acceptance of Phase I capabilities. They
do not require claiming completion of the proposed Phase II research on
collaboration, deception, influence, or human-like reference agents.

\section{Suggested statement of demonstrated capability}

\begin{quote}
Existing experiments demonstrate an EDSL-based multi-agent auction
environment with multiple auction mechanisms, explicit allocation and
payment rules, evolving public market information, and behavioral
comparisons across six LLMs and published human benchmarks. Analysis of 392
completed GPT-4o private-value auctions shows aggregate allocative surplus
efficiency above 99\% in each of five auction settings. Existing rule-based
classifiers identify systematic bidding deviations and differences across
models and mechanisms. The current evidence does not yet demonstrate
operation with ten distinct LLMs; whether the public-information updates
and rule-based classifiers satisfy the corresponding Phase I requirements
depends on their specified definitions.
\end{quote}

\section{Evidence artifacts}
\label{sec:artifacts}

Paths are relative to the repository root. The inspected source revision is
\path{0ba2ecebcfe2f49add765a337dd81aba5c68e89c}; the Phase I audit script,
generated audit outputs, and this memo were added locally after that revision.
The manuscript and data map describe the September 2026 paper revision.
This document is self-contained and does not load external figures, tables,
or a bibliography at compile time.

\small
\begin{longtable}{@{}L{0.06\textwidth}L{0.88\textwidth}@{}}
\toprule
ID & Artifact and evidentiary role \\
\midrule
\endfirsthead
\toprule
ID & Artifact and evidentiary role (continued) \\
\midrule
\endhead
\bottomrule
\endfoot
E1 & \path{writeup/main.tex}; \path{README.md}. Current manuscript and project overview: experimental design, model comparisons, findings, and limitations. \\
E2 & \path{data_map.md}; \path{experiment_logs/V10/}; \path{robustness_logs/}; \path{experiment_logs/paper_interventions/}. Provenance, raw outcomes, configurations, prompts, human reconstruction methods, and missing-data qualifications. \\
E3 & \path{scripts/main.py}; \path{src/util_plan.py}. Experiment orchestration, EDSL model configuration, bidder instantiation, sealed-bid and clock interaction, allocation, payments, and information updates. \\
E4 & \path{src/util_ebay.py}; \path{Prompt/ebay_asking.txt}; \path{scripts/plots/ebay_closing_rule_only.py}; \path{results/analysis/ebay_closing_rule_summary.csv}. Dynamic proxy-auction mechanics, public price histories, and recorded timing comparisons. \\
E5 & \path{results/phase1_evidence/recorded_models.csv}. Inventory of six model identifiers found in recorded CSV exports, including supplementary framing data. \\
E6 & \path{scripts/plots/figure2_models.py}; \path{results/analysis/direction_shares_models.csv}; \path{writeup/figures/figure2_models_direction.png}. Rule-based classification and per-model bid-direction results. \\
E7 & \path{results/analysis/model_smad_results.csv}; \path{results/analysis/difficulty_concordance_models.csv}; \path{results/analysis/interventions_models.csv}; \path{plots/auction_human.csv}; \path{writeup/figures/figure1_bands_gp.png}. Model differences, rankings, framing effects, and human benchmark comparisons. \\
E8 & \path{scripts/audit_phase1_evidence.py}; \path{results/phase1_evidence/allocative_efficiency_summary.csv}; \path{results/phase1_evidence/allocative_efficiency_auctions.csv}; \path{results/phase1_evidence/audit.json}. Reproducible efficiency calculation, auction-level provenance and hashes, CSV cross-checks, and recorded discrepancies. \\
\end{longtable}

\end{document}
